\documentclass[a4paper,11pt]{article}
\usepackage[utf8]{inputenc}
\usepackage[english]{babel}
\usepackage[margin=2.5cm]{geometry}
\usepackage{amsmath}
\usepackage{amssymb}
\usepackage{graphicx}
\usepackage{xcolor}
\usepackage{enumitem}
\usepackage{booktabs}

\title{\textbf{Examination UE -- AIST\\Architectures and system engineering}}
\author{March 2024}
\date{}

\begin{document}

\maketitle

\section*{Terms of the examination}

\begin{itemize}
    \item Course materials of the UE and calculator allowed
    \item Use of personal computer or tablet possible for consulting documents. Use of the Internet, e-mail or any other medium is strictly forbidden.
    \item Non-French speakers can write in English
    \item Parts (A, B) are independent and must be written on separate copies.
\end{itemize}

\section{RF terrestrial and satellite communication links (45 mn)}

The 71-76 GHz and 81-86 GHz bands (E-band) can now be used for very-high-speed communications. Commercial systems operate these frequency bands for point-to-point terrestrial wireless links with data rates of up to 10 Gb/s, using QAM modulation schemes. The E-band is also being considered for links with Low Earth Orbit (LEO) satellites.

We propose to study a terrestrial and satellite radio link, under clear-sky conditions, at 83.5 GHz.

The transmit and receive architecture for the terrestrial link is depicted in Fig. 1 and the specifications of the components are given in Table I.

\begin{figure}[h]
    \centering
    \textit{Fig. 1: transmit (Tx) and receive (Rx) architecture for the terrestrial link}
\end{figure}

\begin{table}[h]
\centering
\caption{Specifications of the components in the terrestrial transmission system}
\textit{[Detailed specifications provided in the exam paper]}
\end{table}

At 83.5 GHz, the lineic attenuation due to the atmosphere is about 0.5 dB/km. The available channel frequency bandwidth for data transmission is $B = 5$ GHz.

\subsection{Free space losses}

\begin{enumerate}[label=\alph*.]
    \item Specify the required conditions for a free-space propagation and give the general expression of the free space loss $A_{FS}$ as a function of distance $d$ and frequency $f$.
    
    \item Calculate $A_{FS}$ at $d = 40$ km (terrestrial link) and $d = 2000$ km (LEO link).
    
    \item Calculate the total free space attenuation $A_T$ (in dB) taking into account the atmosphere.
    
    \item A $10^{-3}$ bit error rate with a QPSK modulation requires a minimum received power at the receiver input of $P_{R,min} = -84$ dBm for the considered bandwidth. Calculate the required EIRP (in dBm) for both the terrestrial and satellite links assuming identical antennas of gain $G = 47$ dBi at the transmitter and receiver.
\end{enumerate}

In the following questions (e, f, g), the Tx and Rx antenna efficiency is $\eta = 0.55$

\begin{enumerate}[label=\alph*., start=5]
    \item In the terrestrial link, the Tx et RX antennas are supposed to be identical circular parabolic antennas. Calculate the required antenna diameter to achieve the gain of $G = 47$ dBi. Is this realistic for a terrestrial point-to-point link?
    
    \item In the satellite link, the ground station is equipped with a parabolic antenna (diameter $d = 1.2$ m). Calculate its gain and deduce the required gain for the satellite antenna.
    
    \item Calculate the 3dB angle (in degrees) of the antennas in the terrestrial link. Comment the need for accurate antenna pointing.
\end{enumerate}

\subsection{Terrestrial link budget}

\begin{enumerate}[label=\alph*., start=8]
    \item What are the Tx and Rx architecture types (see Fig. 1)? Calculate the LO frequency for the Tx and Rx stages.
    
    \item Calculate the noise figure of the receiver (in dB) considering the components between the antenna output and the IF signal output.
    
    \item Calculate the received signal power (in dBm) for the terrestrial link.
    
    \item Calculate the total noise power at the receiver input (in dBm).
    
    \item Calculate the C/N ratio (in dB) and the margin with respect to the required C/N for QPSK modulation.
    
    \item Given the curves in Fig. 2 below, what are the limits of the studied link as dimensioned?
\end{enumerate}

\begin{figure}[h]
    \centering
    \textit{Fig. 2: lineic attenuation (dB/km) due to the rain depending on the frequency}
\end{figure}

\newpage

\section{Maximal link budget of a DWDM link without in-line amplifiers (45 mn)}

One wants to estimate the maximal link budget of a point-to-point DWDM system for Data Center Interconnection purposes in a metropolitan context.

A booster amplifier and a preamplifier are used at the transmitting and receiving terminals respectively but in-line amplifiers are not allowed. Systems must be in conformance with the transponders' characteristics (see annex) and present a 2 dB OSNR margin at the receiver side.

System configuration noted A (resp. B and C) is using only transponder model A (resp. B and C).

\subsection*{Questions}

\begin{enumerate}
    \item Draw a system architecture diagram showing all the optical elements. Each optical element is represented by a rectangle. Patch cords needed to interconnect optical elements are represented by a line connecting rectangles. S and R points have to be represented.
    
    The S point is defined as the entry point of the patch cord to the line at the transmitting side whereas the R point is defined as the output point of the patch cord of the fiber line at the receiving side.
    
    \vspace{4cm}
    
    \item What is the maximum number $N_A$ (resp. $N_B$ and $N_C$) of optical channels that can be transmitted in system configuration A (resp. B and C)?
    
    \vspace{3cm}
    
    \item What is the optical power per channel (in dBm) launched at point S in system configuration A (resp. B and C) when $N_A$ (resp. $N_B$ and $N_C$) channels are transmitted?
    
    \vspace{3cm}
    
    \item What is the OSNR value (in dB) at point S in system configuration A (resp. B and C)? Do you really need to take the booster amplifier's noise into account? Explain why.
    
    \vspace{3cm}
    
    \item What is the minimum power per channel (in dBm) to be received at point R in system configuration A (resp. B and C) if optical channels are received by transponders under power and OSNR threshold conditions with 2dB OSNR margin?
    
    \vspace{3cm}
    
    \item What is the link budget in dB (calculated between S and R points) for system configuration A (resp. B and C)?
    
    \vspace{3cm}
    
    \item What is then the maximum transmission distance for system configuration A (resp. B and C) over a fiber line with 0.18 dB/km attenuation?
    
    \vspace{3cm}
\end{enumerate}

\newpage

\section*{Annex: Card Data Sheet}

Transponder cards comprise a transmitting optical interface and a receiving optical interface to terminate a single optical channel.

\subsection*{100 Gbit/s transponder card, model TRX-A:}

\textit{[Detailed specifications provided in the exam paper]}

\subsection*{200 Gbit/s transponder card, model TRX-B:}

\textit{[Detailed specifications provided in the exam paper]}

\subsection*{400 Gbit/s transponder card, model TRX-C:}

\textit{[Detailed specifications provided in the exam paper]}

\textbf{Note 1:} OSNR is by definition measured in an optical bandwidth $B_{opt}$ of 12.5 GHz.

\subsection*{Optical Amplifier Card:}

\textit{[Detailed specifications provided in the exam paper]}

\textbf{Note 1 :} The gain figure is given under small input signal conditions when the amplifier is not saturated.

The amplified spontaneous emission (ASE) power at the output of an optical amplifier with gain $G$ and noise factor $n_{sp}$ (both in linear scale) is given by:

$$P_{ASE} = 2n_{sp}h\nu(G-1)B_{opt}$$

For a C-band central wavelength of 1545 nm with $B_{opt}=12.5$ GHz, this formula becomes:

$$P_{ASE,dBm} = -58 + F_{dB} + G_{dB}$$

\subsection*{Fixed grid C-band 40 channels MUX/DEMUX card}

This card comprises 40 ports for optical channels and one port for WDM multiplex. This card is suited for IUT-T 100 GHz fixed grid WDM.

\subsection*{Fixed grid C-band 80 channels MUX/DEMUX card}

This card comprises 80 ports for optical channels and one port for WDM multiplex. This card is suited for IUT-T 50 GHz fixed grid WDM.

\end{document}
